\graphicspath{{./},{./intro_figs}}
\section{\texorpdfstring{\MakeUppercase{Introduction}}{Introduction}}
\label{sec:intro}

\counterwithin{figure}{section}

\subsection{Multimessenger Probes of AGN Disks}

\subsection{Introduction}
\label{sec:intro-ch2}

Binary black hole (BBH) mergers are observed in gravitationals waves (GW) at an increasing rate as LIGO-Virgo-KAGRA (LVK) improve detector sensitivities \citep[e.g.][ see also \footnote{https://gracedb.ligo.org/superevents/public/O4/}]{Abbott+20-GW190521astro}.
Observed BBH include black holes (BH) with individual masses in the expected upper mass gap ($\sim 50-120M_{\odot}$) where BHs are not expected from stellar evolution \citep{Belcynzski+16,RenzoSmith24} due to the pulsational pair instability.
A natural explanation for such BHs is via a series of hierarchical mergers that may occur in a dynamical environment, such as globular clusters \citep{Rodriguez+18}, nuclear star clusters (NSC) \citep{Antonini&Rasio16,Hoang+18,Fragione+19} and active galactic nuclei (AGN) accretion disks \citep{McKernan+12,McKernan+14,Bartos+17,StoneMetzgerHaiman17, ArcaSedda23,FordMcKernan25}.
Among hierarchical merger models, AGN are notable for both the efficiency of BBH formation as well as for retaining most post-merger products (as even strong kicks from mergers correspond to small orbital perturbations in such a deep potential well).
Action of gas as a dynamical coolant and orbital migration producing relatively low-velocity dynamical encounters drive BBH formation \citep[e.g.][]{Rowan+23,QianLiLai24,DodiciTremaine24}.
The result is a high rate of intermediate mass black hole (IMBH) formation relative to other channels \citep{McKernan+12,Bellovary+16,Secunda+21,Yang+19,McKernan+20a,TagawaHaimanKocsis20,Vaccaro+24}, though if gap formation by IMBHs occurs, their presence may disrupt migration traps that would otherwise facilitate high rates of hierarchical mergers \citep{Gilbaum+25} but could also introduce additional traps at the gap edges \citep{Horn+12}.
The rate of BBH mergers from AGN is also expected to be significantly higher than from quiescent NSCs \citep{FordMcKernan22}.

The observed population of BBH mergers intrinsically favors positive effective spin \citep{Abbott+23-O3b}, implying formation channels bias spin directions towards alignment with BBH orbital angular momentum.
This is not expected from gas-free hierarchical merger channels, but can occur in the AGN channel as BH are torqued towards alignment with the gas disk over time \citep{Tagawa+20spin,Vajpeyi+22,McKernanFord24}.
Additionally, there is a possible anti-correlation observed in the mass ratio $q$ and effective spin $\Xeff$ distribution of BBH mergers \citep{Callister+21}.

This seems difficult to achieve in standard variants of other dynamical merger channels (e.g. globular clusters and NSC although see \citet{Su25} for a discussion of how such an anti-correlation might arise in a dynamical triple channel).
In the isolated binary channel, \citep[e.g.][]{Olejak+24,BanerjeeOlejak24} demonstrate the possibility of generating a ($q,\Xeff$) anti-correlation via stable mass transfer.
In the AGN channel, the $(q,\Xeff)$ can arise due to several biases \citep{Mckernan+22:q-Xeff,Santini+23, DittmannDempseyLi24}.
\citet{Vaccaro+24} show that if gas hardening is efficient in AGN, then the $(q,\Xeff)$ distribution should extend to significantly smaller values of $q$ and high $\Xeff$.
Thus, LVK O4 results can directly test gas hardening efficiency in AGN channel models.

Note that the AGN channel accounts for a still uncertain fraction ($f_{\rm AGN,BBH}$) of the observed LVK GW events.
We will learn considerably about the average properties of AGN and NSCs out to $z \sim 1$ by constraining $f_{\rm AGN, BBH}$, \emph{regardless of its actual value}, see \citet{FordMcKernan25} for a more detailed discussion of this point.
Thus, in presenting distributions of AGN channel events in ($q,\Xeff$) as a function of model choices, we are not claiming that AGN are responsible for all (or any) of the observed events.
Rather we are demonstrating areas of AGN parameter space that might be ruled out by comparing with LVK O4 (O5) results.
This in turn hopefully allows us to constrain $f_{\rm AGN,BBH}$, the fraction of mergers originating from AGN.

%Here we present a study using the new \code{McFACTS} code \citep[Monte Carlo For AGN Channel Testing and Simulations:][hereafter Paper I]{McKernan+25-mcfactsI} to test the predictions of the AGN channel for the $q$--$\Xeff$ distribution across a wide range of parameter space in both AGN disk and NSC models, with a view to tightly constraining requirements on the AGN channel.
%\citet{Delfavero+25} presents the first study of the AGN channel in a semi-realistic Universe of AGN and NSCs.
%\code{McFACTS} features substantial improvements to the Monte Carlo code used in \citet{McKernan+20a} to explore $\Xeff$ of mergers in AGN and \citet{Mckernan+22:q-Xeff} who explored the $q$--$\Xeff$ anti-correlation.
%The greater flexibility for exploring parameter space includes many improved models for disk-NSC interactions such as dynamical encounters, orbital migration, and binary evolution.



\subsection{Introduction} 
\label{sec:intro-ch3}

Galactic nuclei are typically composed of a SMBH and a nuclear star cluster \citep{Ghez+98,Ghez+08}.
When gas accretes onto the SMBH, some of the objects in the now active galactic nucleus will be embedded in the accretion disk while others will have disk-crossing orbits.
A hypothesis has been suggested that the disk may generate chains of hierarchical mergers between stellar-mass black holes that produce intermediate-mass remnants \citep{McKernan+14, Bartos+17, Stone+17, FordMcKernan25, Abbott+20-GW190521, LVK+25-GW231123}, much like the collisional growth model for planet formation in circum-stellar disks \citep{Pollack+96}.
These mergers could produce measurable electromagnetic counterparts from interactions with the AGN gas which would be otherwise absent for two black holes merging in the field \citep{McKernan+19}.
Yet, other energetic events such as core-collapse SNe could become a source of noise that must be accounted for in the search for these merger counterparts \citep{Graham+20}.

In addition to those stars initially embedded in the disk when it forms, self-gravitating AGN disks can naturally fragment and form new stars, particularly at large radii \citep{ShlosmanBegelman89, CollinZahn08}, though magnetic fields \citep{Gerling-Dunsmore+25, Tsung+25} and feedback from embedded objects \citep{HanklaJiangArmitage20} may inhibit collapse.
Furthermore, AGN disks can quickly capture stars from the inclined population \citep{Fabj+20,MacLeodLin20}.
Any massive embedded stars may detonate as SNe within the disk confines, though stellar evolution in AGN disks may proceed very differently than in the interstellar medium \citep[ISM,][]{CantielloJermynLin21}.
Even in this case, SNe may yet result from stellar evolution \citep{Jermyn+21} or collisions.
SNe in the dusty outskirts of AGN disks may have already been detected as bright optical and infrared flares \citep[e.g.][]{Assef+18}.

We are motivated to explore the evolution of SNe embedded in an AGN disk due to the possibility that AGN flares could be electromagnetic counterparts to the gravitational wave detections from binary black hole (BBH) merger events, one particular case to underscore being GW190521 \citet{Abbott+20-GW190521}.
While the observed flare \citep{Graham+20} did not resemble a standard SN light curve, embedded SNe could behave differently than SNe occurring in the ISM. 
Thus, by characterizing SN light curves in AGN, we can rule out false positive counterparts to BBH mergers as well as constrain the rate of SNe in AGN. 
Looking forward, detections of embedded SNe (or other transients in AGN disks) will allow us to constrain the composition, evolution, and dynamics of nuclear star clusters that coexist with an AGN disk, as well as to probe the disk properties.

Whether the radiation from SNe escapes from an AGN disk is a separate matter from the mechanical breakout of the shockwave.
The non-uniform density and variable disk thickness have large impacts on a SN shock's morphological in comparison to their galactic counterparts \citep{MacLowMcCray88}.
The right combination may slow down a SN shock significantly or prevent it from escaping entirely, thus depositing their kinetic energy in the disk.
These muffled SNe may not punch holes through the disk but could still be a contributing energy source for driving turbulence and maintaining pressure support in disks that would otherwise collapse under self-gravity \citep{Rozyczka+95,Moranchel+21}.

\citet{Grishin+21} performed two dimensional (2D) hydrodynamic simulations in the $r$-$z$ plane of SNe occurring in a flat slab of gas at the midplane and with a vertical offset and find the shock breaks free from the disk above and below the midplane for both cases.
\citet{Moranchel+21} ran three-dimensional simulations in spherical coordinates for the upper half of the disk and find agreement that shocks from SNe at the midplane broke free from the disk.
We expand upon this work by estimating where SNe may break free from two AGN disk models and perform three-dimensional hydrodynamic simulations of a SN embedded in a vertically stratified shearing disk for midplane and offset SNe explosions.

Galactic nuclei contain stars, some of which, if embedded in an active galactic nucleus (AGN) disk could detonate as embedded supernovae (SNe).
Self-gravitating AGN disks can naturally fragment to form stars particularly at large radii \citep[e.g.][]{Shlosman89,Collin08}, which could evolve to SNe.
Nuclear cluster stars can be quickly captured into AGN disks  \citep{Fabj20,MacLeod20}, and can also evolve to SNe.
Stellar evolution in AGN disks may proceed very differently than in vacuum \citep{Cantiello+21}, but even in this case SNe may yet result from stellar evolution \citep{Jermyn21} or collisions, and contribute to disk turbulence \citep{Moranchel+21}.
SNe in the dusty outskirts of AGN disks may have already been detected as bright optical and infra-red flares \citep[e.g.][]{Assef18}.

AGN disks are receiving renewed interest as potential sites for stellar mass black hole (BH) mergers. BHs embedded in AGN disks experience gas torques and can migrate \citep[e.g.][]{Syer91,Levin07,McKernan+12}. Differential migration leads to encounters with other BH (as well as neutron stars and stars) and binary formation. Binaries in AGN can be hardened or softened via dynamical encounters \citep[e.g.][]{Leigh+18,Samsing20,Wang21} or gas torques  \citep[e.g.][]{Tiede20,Derdzinski21,YapingLi21}. Thus we expect mergers in AGN disks detectable with gravitational waves \citep{McKernan+14,Bartos+17,Stone+17,Secunda+19} and potentially detectable electromagnetically \citep{McKernan+19,Perna21}. Such mergers can occur at migration traps \citep[e.g.][]{Lyra+10,Sandor+11,Horn+12,Bellovary2016,Yang+19,Secunda+20}, in the bulk disk \citep[e.g.][]{Tagawa20a,McK20a} and kicked outside the disk itself \citep{Tagawa20b}.

We are particularly motivated to consider the light curve of a SN embedded in an AGN disk by the first candidate electromagnetic counterpart to a gravitational wave detected binary black hole (BBH) merger.
\citet{Graham+20} reported a flare in AGN J124942.3+344929 that was observed by the Zwicky Transient Facility (ZTF) and designated ZTF19abanrhr.
Since the flare was inside the source volume of the gravitational wave detection GW190521 \citep{Abbott+20-GW190521}
the authors suggest it could have been produced by the recoiling remnant black hole if the merger had occurred within the AGN accretion disk and experienced ram-pressure stripping of the gas bound within its Hill sphere by the disk \citep{McKernan+19}.
Using this model \citet{Graham+20} predicted a black hole mass $M\approx150\,\Msun$ which was later confirmed by the LIGO Scientific Collaboration and the Virgo Collaboration \citep{Abbott+20-GW190521}.


They dismissed a SN causing the flare because it did not show any color evolution as is typical of SNe light curves \citep{Foley+11}.
However, the integrated energy of the flare was $\sim10^{51}$ erg s$^{-1}$, which is conspicuously the typical energy budget of a SN.
One class of SNe are denoted as superluminous SNe (SLSNe) since their light curves peak above $\sim 10^{44}\,\text{erg s}^{-1}$ \citep{Gal-Yam12}.
One possible origin of these SLSNe is from stars that experienced extensive mass loss during their final years pre-SN \citep{Moriya+11}.
The amount of radiated energy increases with the enevelope mass until the mass is similar to the ejecta mass, and the width of the light curve increases with the envelope size \citep{GinzburgBalberg12}.
If embedded SNe interact with a dense AGN disk in a similar way, their radiation may likely also be larger than their ``naked" cousins whereby the disk reprocesses much of their mechanical energy into radiation as the shock heats the disk.

Characterizing light curves and color evolution from these embedded SNe is crucial for accurately distinguishing between potential transient events present in AGN disks during archival searches and live rapid-response follow-up campaigns for counterparts to gravitational wave events.
Furthermore, by understanding how SNe appear in AGN, we can rule out false positive counterparts to black hole mergers as well as constrain the rate of SNe n AGN.
Looking forward, detections of embedded SNe would allow us to constrain the stellar populations, evolution, and dynamics of nuclear star clusters (NSCs) that coexist with an AGN disk, as well as the properties of the disk itself.

